\chapter{Finite volume method (FVM)}
\label{chap:dev:finite-volume-method}

\begin{quote}
Phase metadata. Spec dir:
\passthrough{\lstinline!specs/2026-05-04-phase-14-finite-volume-method/!}.
Inserted into the roadmap on 2026-05-04 between the just-merged
Phase 13 (Vector and tensor fields) and the originally-numbered
Phase 14 (Functional evaluation for vector/tensor data, now Phase 15);
all subsequent phases renumbered down by one.
Adds the conservative cell-flux discretization
\(\nabla\cdot(D\,\nabla\psi)\) for heterogeneous diffusivity. Version
target: \passthrough{\lstinline!0.14.0!}.
\end{quote}

\subsection{1. Theory}\label{phase14-theory}

\subsubsection{Conservation laws and the flux-form
PDE}\label{phase14-conservation}

The heterogeneous-coefficient diffusion equation
\[
\partial_t \psi(\mathbf x, t) \;=\; \nabla\cdot\!\big[D(\mathbf x)\,\nabla\psi(\mathbf x, t)\big]
\]
is a conservation law for the scalar density \(\psi\):
\(\partial_t \psi + \nabla\cdot\mathbf J = 0\) with the diffusive flux
\(\mathbf J(\mathbf x) = -D(\mathbf x)\,\nabla\psi(\mathbf x)\). Integrating over
any closed control volume \(V\) gives
\[
\frac{d}{dt}\int_V \psi\,dV \;=\; -\oint_{\partial V} \mathbf J \cdot d\mathbf S,
\]
i.e.\ the rate of change of \(\int_V \psi\) equals the net flux
through \(\partial V\). On a periodic domain, every face is shared
between two cells with equal-and-opposite contributions, so the
total mass \(\int \psi\,dV\) is exactly conserved. The
finite-\emph{volume} method discretizes this integrated form
directly and inherits the conservation property at every step
without further assumptions; the finite-\emph{difference} method
discretizes the differential form and conserves only when the
underlying stencil happens to be telescoping (which it is, for the
3-point central-difference Laplacian on a uniform grid). For
heterogeneous \(D\) the FVM cell-flux formulation generalizes
naturally; the FD formulation of \(\nabla\cdot(D\nabla\psi)\) by
operator composition would still converge but the conservation and
flux-matching properties are not built-in.

\subsubsection{Why harmonic-mean face
\(D\)}\label{phase14-harmonic}

Consider a steady-state 1D problem on two adjacent cells with
\(D = D_L\) on the left and \(D = D_R\) on the right, with a known
flux \(J\) traversing the interface. The continuous
heat-equation steady state has constant flux
\(J = -D(x)\,d\psi/dx\). Integrating across the cell-pair:
\[
\psi_R - \psi_L \;=\; \int_{x_L}^{x_R} \frac{d\psi}{dx}\,dx \;=\; -\,J\!\int_{x_L}^{x_R}\!\frac{dx}{D(x)}.
\]
For the piecewise-constant \(D\), the integral evaluates to
\(\frac{h_L/2}{D_L} + \frac{h_R/2}{D_R}\) on a uniform grid with
cell size \(h\):
\[
\psi_R - \psi_L \;=\; -\,J\,\frac{h}{2}\Big(\frac{1}{D_L} + \frac{1}{D_R}\Big).
\]
Setting this equal to the FVM update \(\psi_R - \psi_L = -\,J\,h\,/\,D_{\text{face}}\)
(with whatever face-centered \(D_{\text{face}}\) we choose) and
solving:
\[
D_{\text{face}} \;=\; \frac{2\,D_L\,D_R}{D_L + D_R}.
\]
This is the harmonic mean. The arithmetic mean
\(\tfrac12(D_L + D_R)\) over-predicts the steady-state flux
whenever the two values differ; the harmonic mean recovers it
exactly. The Phase 14 implementation uses the harmonic mean at
every face --- LeVeque (2007) and Patankar (1980) both make this
the textbook choice for diffusion under non-uniform
coefficients.

\subsection{2. Math derivation}\label{phase14-math}

\subsubsection{Cell-flux discretization}\label{phase14-cell-flux}

For cell \((i, j, k)\) with size \(h_a\) along axis \(a\) and the
six neighboring cells along the three axes, write the time
derivative as the negative divergence of the flux:
\[
\partial_t\psi(i, j, k) \;=\; \frac{1}{h_x h_y h_z}\bigg[\,\sum_{a \in \{x, y, z\}}\!\big(F^{(a,+)}_{ijk} - F^{(a,-)}_{ijk}\big)\,A_a\bigg],
\]
where \(A_a = h_b h_c\) is the face area perpendicular to axis
\(a\). The flux through each face is approximated by a
central-difference of \(\psi\) across that face, scaled by the
face-centered \(D\):
\[
F^{(x,+)}_{ijk} \;=\; -\,D^{(x,+)}_{ijk}\,\frac{\psi(i+1, j, k) - \psi(i, j, k)}{h_x},
\]
and similarly for the other 5 faces. With
\(D^{(x,+)}_{ijk} = \tfrac{2\,D(i,j,k)\,D(i+1,j,k)}{D(i,j,k) + D(i+1,j,k)}\)
and the cell volume \(h_x h_y h_z\) cancelling against \(A_a / h_a\),
the per-cell time derivative is
\[
\partial_t\psi(i,j,k) \;=\; \sum_a \frac{1}{h_a^2}\,\Big[\,D^{(a,+)}_{ijk}\,(\psi_{a,+} - \psi)\,-\,D^{(a,-)}_{ijk}\,(\psi - \psi_{a,-})\,\Big].
\]
This is precisely the expression
\passthrough{\lstinline!fvm::cellLaplacian!} computes per cell.

\subsubsection{Order-2 truncation analysis}\label{phase14-order-2}

For smooth \(D\), \(\psi\) and uniform spacing \(h\), expand each
neighbor-value about the cell center:
\[
\psi(i \pm 1) \;=\; \psi \pm h\,\partial_x\psi + \tfrac{h^2}{2}\,\partial_x^2\psi \pm \tfrac{h^3}{6}\,\partial_x^3\psi + O(h^4),
\]
\[
D(i \pm 1) \;=\; D \pm h\,\partial_x D + \tfrac{h^2}{2}\,\partial_x^2 D + O(h^3),
\]
\[
D^{(x,\pm)} \;=\; D \pm \tfrac{h}{2}\,\partial_x D + \tfrac{h^2}{8}\,\partial_x^2 D + O(h^3) \pm O(h\,\partial_x D \cdot D / D^2) \dots
\]
Combining the harmonic-mean expansions of \(D^{(x,\pm)}\) with the
finite-difference \(\psi\) increments, the cell-flux expression
collapses to
\[
\partial_t\psi(i,j,k) \;=\; \nabla\!\cdot\!(D\,\nabla\psi) \;+\; O(h^2),
\]
matching the continuous operator at second order. The
order-2 convergence rate is verified empirically by the test fixture
\passthrough{\lstinline!AnalyticalSolutionConvergence\_Order2!} on
the manufactured solution
\(D(x) = 1 + 0.5\cos(x)\), \(\psi(x) = \sin(x)\) (closed-form
\(\nabla\cdot(D\nabla\psi) = -\sin(x) - 0.5\sin(2x)\)) and by
\passthrough{\lstinline!RecoversAnalyticalEigenfunctionAtOrder2!}
on the standard \(\sin\!\cdot\!\sin\!\cdot\!\sin\) eigenfunction
under uniform \(D\).

\subsubsection{Generalized CFL bound}\label{phase14-cfl}

The continuous heat-equation eigenvalue spectrum of the
\(-\nabla\!\cdot\!(D\nabla\,\cdot)\) operator is bounded above by
\(D_{\max}\!\cdot\!\lambda_{\max}^{(\text{Laplacian})}\). On a
uniform Cartesian grid with the FVM cell-flux discretization, the
discrete eigenvalues are bounded by
\(D_{\max}\!\cdot\!\sum_a (4/h_a^2)\) (each axis contributes the
range \([-4/h_a^2, 0]\) to the FD-Laplacian spectrum, and the
harmonic-mean coefficients are bounded by \(D_{\max}\) above).
Forward-Euler stability requires
\(D_{\max}\!\cdot\!dt\!\cdot\!\sum_a(1/h_a^2)\,\le\,0.5\); RK4
admits 0.6963 by the same construction. The
\passthrough{\lstinline!solve::diffuse(psi, D\_field, dt, n, tag)!}
overload reuses
\passthrough{\lstinline!detail::checkDiffusionCFL<Tag>!} from
Phase 6 with \(D = D_{\max}\) substituted for the scalar
diffusivity.

\subsection{3. Algorithm}\label{phase14-algorithm}

\subsubsection{File layout}\label{phase14-files}

\begin{itemize}
\tightlist
\item
  \href{../../../include/gridcalc/fvm/CellLaplacian.hpp}{\passthrough{\lstinline!include/gridcalc/fvm/CellLaplacian.hpp!}}
  --- new header. Exposes
  \passthrough{\lstinline!fvm::cellLaplacian(psi, D)!}; private
  \passthrough{\lstinline!detail::harmonicMean(a, b)!}.
\item
  \href{../../../include/gridcalc/solve/Diffusion.hpp}{\passthrough{\lstinline!include/gridcalc/solve/Diffusion.hpp!}}
  --- extended with the heterogeneous-\(D\) overload. The constant-\(D\)
  Phase 5 / Phase 6 overload is unchanged.
\item
  \href{../../../examples/heterogeneous_diffusion.cpp}{\passthrough{\lstinline!examples/heterogeneous\_diffusion.cpp!}}
  --- demo binary; CLI for grid size, time step, integrator, D
  parameters, IC parameters, output dir.
\item
  \href{../../../examples/common/heterogeneous_diffusion.hpp}{\passthrough{\lstinline!examples/common/heterogeneous\_diffusion.hpp!}}
  --- shared helpers
  (\passthrough{\lstinline!makeGaussianIc!},
  \passthrough{\lstinline!makeHeterogeneousD!},
  diagnostics).
\item
  \href{../../../test/fvm_cell_laplacian_test.cpp}{\passthrough{\lstinline!test/fvm\_cell\_laplacian\_test.cpp!}},
  \href{../../../test/diffusion_test.cpp}{\passthrough{\lstinline!test/diffusion\_test.cpp!}}
  (heterogeneous-D additions),
  \href{../../../test/fvm_diffusion_acceptance_test.cpp}{\passthrough{\lstinline!test/fvm\_diffusion\_acceptance\_test.cpp!}},
  \href{../../../test/heterogeneous_diffusion_test.cpp}{\passthrough{\lstinline!test/heterogeneous\_diffusion\_test.cpp!}}
  --- 14 new tests across four files.
\item
  \href{../../../scripts/render_diffusion_snapshots.py}{\passthrough{\lstinline!scripts/render\_diffusion\_snapshots.py!}}
  --- matplotlib helper that reads the demo's VTK snapshots and
  renders \(z\)-midplane PNGs for the User Guide chapter.
\end{itemize}

\subsubsection{The \texttt{cellLaplacian}
loop}\label{phase14-loop}

The hot path is a triple-nested loop with six harmonic-mean calls
plus six \(\psi\)-difference / scaled additions per cell. No
intermediate field is materialized; the result is written directly
into the output \passthrough{\lstinline!Field<double>!}:

\begin{lstlisting}[language={C++}]
for (int k = 0; k < grid.getNz(); ++k)
  for (int j = 0; j < grid.getNy(); ++j)
    for (int i = 0; i < grid.getNx(); ++i) {
      const double psi_c = psi(i, j, k);
      const double D_c   = D(i, j, k);

      const double D_face_xp = detail::harmonicMean(D_c, D(i + 1, j, k));
      const double D_face_xm = detail::harmonicMean(D(i - 1, j, k), D_c);
      // ... y, z faces ...

      const double rhs_x = (D_face_xp * (psi(i + 1, j, k) - psi_c)
                          - D_face_xm * (psi_c - psi(i - 1, j, k))) * inv_hx2;
      // ... y, z axes ...
      result(i, j, k) = rhs_x + rhs_y + rhs_z;
    }
\end{lstlisting}

\subsubsection{Heterogeneous-\(D\)
\passthrough{\lstinline!solve::diffuse!}}\label{phase14-solver}

The new overload validates the harmonic-mean contract
(\(D > 0\) everywhere) and computes \(D_{\max}\) for the CFL bound
in a single pre-loop pass:

\begin{lstlisting}[language={C++}]
double D_max = 0.0;
for (std::size_t i = 0; i < D.getSize(); ++i) {
  if (!(D.data()[i] > 0.0)) throw std::invalid_argument("...");
  if (D.data()[i] > D_max) D_max = D.data()[i];
}
detail::checkDiffusionCFL<Integrator>(psi.getGrid(), D_max, dt);
integrate(psi,
          [&D](const Field<double>& f) { return fvm::cellLaplacian(f, D); },
          dt, n_steps, tag);
\end{lstlisting}

The closure
\passthrough{\lstinline![\&D](const Field<double>\& f) \{ ... \}!}
captures \(D\) by reference; its lifetime extends through the
integrator call, which is the entire scope of
\passthrough{\lstinline!solve::diffuse!}. No \(D\) copy is made.

\subsubsection{Demo + CI gate}\label{phase14-demo}

The demo binary
\passthrough{\lstinline!examples/heterogeneous\_diffusion.cpp!}
threads CLI options through
\passthrough{\lstinline!parseOptions!}, builds the IC and the
\(D\)-field via the
\passthrough{\lstinline!examples/common/!} helpers, and time-steps in
chunks of \passthrough{\lstinline!--snapshot-every!} steps writing
VTK snapshots between chunks. The CI gate
\passthrough{\lstinline!test/heterogeneous\_diffusion\_test.cpp!}
runs a stripped-down version (\(N=24\), 4 \(\times\) 50 RK4 steps,
~1.6 s on Apple Clang debug) once per process via a static-cached
lambda; three TEST\_F bodies share the cached snapshots.

\subsection{4. Design decisions}\label{phase14-decisions}

These distill the choices recorded in
\passthrough{\lstinline!specs/2026-05-04-phase-14-finite-volume-method/requirements.md!}.

\begin{itemize}
\item
  \textbf{Hybrid incorporation, not replacement.} New
  \passthrough{\lstinline!gridcalc::fvm!} layer hosting the
  cell-flux Laplacian and its diffusion-solver dispatch. Phase 1's
  \passthrough{\lstinline!diff::laplacian!} and Phase 5 / Phase 6's
  constant-\(D\)
  \passthrough{\lstinline!solve::diffuse!} are unchanged. On uniform
  grids with constant \(D\), the FVM cell-flux discretization
  produces the same stencil numbers as the FD path; the new code
  is only physically distinguishable from the old when
  \(D\) varies.
\item
  \textbf{Harmonic-mean face \(D\), not arithmetic.} Recovers the
  correct steady-state flux across \(D\)-discontinuities (see
  \Cref{phase14-harmonic}). The arithmetic mean would give a
  superficially valid order-2 convergent scheme on smooth \(D\) but
  the wrong leading-error constant on \(D\) jumps. Documented in
  the public Doxygen block on
  \passthrough{\lstinline!fvm::cellLaplacian!}.
\item
  \textbf{Reuse the existing \passthrough{\lstinline!solve::*!}
  integrator tags.} No new \passthrough{\lstinline!fvm::!}-prefixed
  tag set. The new
  \passthrough{\lstinline!solve::diffuse(psi, D\_field, ...)!}
  overload accepts
  \passthrough{\lstinline!ExplicitEuler\{\}!} and
  \passthrough{\lstinline!RK4\{\}!} by design --- one less surface
  to maintain, and the heterogeneous path inherits Phase 6's
  CFL-bound machinery via
  \passthrough{\lstinline!detail::checkDiffusionCFL<Tag>!}.
\item
  \textbf{Periodic-only at Phase 14.} Phase 19 (the renumbered
  non-periodic-BC phase) handles Dirichlet / Neumann walls; FVM is
  the natural foundation for that work because the flux
  formulation makes wall conditions drop out as boundary-flux
  prescriptions rather than ghost-cell constructions.
\item
  \textbf{No higher-order FVM.} 2nd-order central-difference face
  fluxes only. Higher-order MUSCL / WENO reconstructions are out
  of scope.
\item
  \textbf{The textbook demo's manufactured-solution trap.}
  Initially we tried the naive manufactured solution
  \(\psi(x, y, z, t) = \exp(-D(x)\,t)\,\sin(z)\) under
  \(D = D(x)\). It satisfies the PDE at \(t = 0\) but \emph{not} for
  \(t > 0\) --- once \(\psi\) acquires \(x\)-dependence from the
  position-dependent decay rate, the \(x\)-axis flux
  \(\partial_x[D\,\partial_x\psi]\) becomes non-zero and the
  analytical reference drifts away from the discrete trajectory.
  The
  \passthrough{\lstinline!fvm\_diffusion\_acceptance\_test.cpp!}
  source comment records this trap so a future contributor does
  not re-step it. The test file ships two complementary checks
  instead: order-2 convergence on the
  \(\sin\!\cdot\!\sin\!\cdot\!\sin\) eigenfunction with \(D\)
  uniform (analytical reference, valid because \(\psi\) and \(D\)
  remain uncoupled) plus qualitative properties on a genuinely
  heterogeneous \(D\) (mass conservation, max-\(\psi\) monotone
  decrease, positivity preservation).
\end{itemize}

\subsection{5. References}\label{phase14-references}

External:

\begin{itemize}
\tightlist
\item
  \textbf{Patankar, S. V. (1980). \emph{Numerical Heat Transfer and
  Fluid Flow}. Hemisphere Publishing.} ISBN 978-0-89116-522-4.
  Section 4.2.3 derives the harmonic-mean face conductivity used in
  this phase from the requirement that the steady-state flux match
  the continuous answer across cell faces with discontinuous
  diffusivity.
\item
  \textbf{LeVeque, R. J. (2007). \emph{Finite Difference Methods for
  Ordinary and Partial Differential Equations: Steady-State and
  Time-Dependent Problems}. SIAM.}
  \url{https://doi.org/10.1137/1.9780898717839}. Chapter 4
  (``Variable coefficients and conservation laws'') covers the FVM
  cell-flux discretization, the order-2 truncation argument used
  in
  \Cref{phase14-order-2}, and the harmonic-mean / arithmetic-mean
  trade-off; Chapter 9 covers the generalized CFL bound used in
  \Cref{phase14-cfl}.
\item
  \textbf{Hundsdorfer, W., \& Verwer, J. G. (2003).
  \emph{Numerical Solution of Time-Dependent Advection-Diffusion-Reaction
  Equations}. Springer.}
  \url{https://doi.org/10.1007/978-3-662-09017-6}. Chapter 1
  (``Spatial discretization'') covers the FVM treatment of
  heterogeneous diffusion and the conservation-form proof that
  carries through to the time integrator. Used informally in
  \Cref{phase14-conservation}.
\item
  \textbf{VTK File Formats reference.} Permanent URL:
  \url{https://kitware.github.io/vtk-examples/site/VTKFileFormats/}.
  Documents the legacy ASCII
  \passthrough{\lstinline!STRUCTURED\_POINTS!} format the demo
  writes; reused from Phase 12 with no changes.
\end{itemize}

Internal cross-references (do not satisfy the ``external'' minimum):

\begin{itemize}
\tightlist
\item
  \Cref{chap:dev:periodic-3d-scalar-laplacian} --- Phase 1's
  \passthrough{\lstinline!diff::laplacian!}, the FD reference the
  Phase 14 constant-\(D\) agreement gate compares against.
\item
  \Cref{chap:dev:explicit-euler-diffusion} and
  \Cref{chap:dev:rk4--generic-time-integrator} --- Phase 5 / Phase
  6's
  \passthrough{\lstinline!solve::diffuse!} +
  \passthrough{\lstinline!solve::integrate!} dispatch reused by the
  new heterogeneous-\(D\) overload.
\item
  \Cref{chap:dev:cahn-hilliard} --- Phase 12's CH demo + VTK writer
  + acceptance-test pattern, which Phase 14's demo + CI gate
  mirror.
\end{itemize}
